\documentclass[10pt,a4paper]{amsart}
\usepackage{amsmath,amssymb,amsthm,graphicx}
% Retained minimal preamble for the Stacks Project chapter flat.tex (More on Flatness).
% Verbatim excerpt of the chapter preamble: Xy-pic setup (native preamble lines 14-18) and the native
% theorem declarations with the mathematical macros (native preamble lines 154-217). Copied mechanically;
% nothing added, removed or redefined. The wrapper class and the amsmath/amssymb/amsthm/graphicx base
% packages are supplied by the assembler, so the class/externaldocument part of the native preamble is omitted.
% For commutative diagrams we use Xy-pic
\usepackage[all]{xy}

% We use 2cell for 2-commutative diagrams.
\xyoption{2cell}
% Theorem environments.
%
\theoremstyle{plain}
\newtheorem{theorem}[subsection]{Theorem}
\newtheorem{proposition}[subsection]{Proposition}
\newtheorem{lemma}[subsection]{Lemma}

\theoremstyle{definition}
\newtheorem{definition}[subsection]{Definition}
\newtheorem{example}[subsection]{Example}
\newtheorem{exercise}[subsection]{Exercise}
\newtheorem{situation}[subsection]{Situation}

\theoremstyle{remark}
\newtheorem{remark}[subsection]{Remark}
\newtheorem{remarks}[subsection]{Remarks}

\numberwithin{equation}{subsection}

% Macros
%
\def\lim{\mathop{\mathrm{lim}}\nolimits}
\def\colim{\mathop{\mathrm{colim}}\nolimits}
\def\Spec{\mathop{\mathrm{Spec}}}
\def\Hom{\mathop{\mathrm{Hom}}\nolimits}
\def\Ext{\mathop{\mathrm{Ext}}\nolimits}
\def\SheafHom{\mathop{\mathcal{H}\!\mathit{om}}\nolimits}
\def\SheafExt{\mathop{\mathcal{E}\!\mathit{xt}}\nolimits}
\def\Sch{\mathit{Sch}}
\def\Mor{\mathop{\mathrm{Mor}}\nolimits}
\def\Ob{\mathop{\mathrm{Ob}}\nolimits}
\def\Sh{\mathop{\mathit{Sh}}\nolimits}
\def\NL{\mathop{N\!L}\nolimits}
\def\CH{\mathop{\mathrm{CH}}\nolimits}
\def\proetale{{pro\text{-}\acute{e}tale}}
\def\etale{{\acute{e}tale}}
\def\QCoh{\mathit{QCoh}}
\def\Ker{\mathop{\mathrm{Ker}}}
\def\Im{\mathop{\mathrm{Im}}}
\def\Coker{\mathop{\mathrm{Coker}}}
\def\Coim{\mathop{\mathrm{Coim}}}

% Boxtimes
%
\DeclareMathSymbol{\boxtimes}{\mathbin}{AMSa}{"02}

%
% Macros for moduli stacks/spaces
%
\def\QCohstack{\mathcal{QC}\!\mathit{oh}}
\def\Cohstack{\mathcal{C}\!\mathit{oh}}
\def\Spacesstack{\mathcal{S}\!\mathit{paces}}
\def\Quotfunctor{\mathrm{Quot}}
\def\Hilbfunctor{\mathrm{Hilb}}
\def\Curvesstack{\mathcal{C}\!\mathit{urves}}
\def\Polarizedstack{\mathcal{P}\!\mathit{olarized}}
\def\Complexesstack{\mathcal{C}\!\mathit{omplexes}}
% \Pic is the operator that assigns to X its picard group, usage \Pic(X)
% \Picardstack_{X/B} denotes the Picard stack of X over B
% \Picardfunctor_{X/B} denotes the Picard functor of X over B
\def\Pic{\mathop{\mathrm{Pic}}\nolimits}
\def\Picardstack{\mathcal{P}\!\mathit{ic}}
\def\Picardfunctor{\mathrm{Pic}}
\def\Deformationcategory{\mathcal{D}\!\mathit{ef}}

% --- appended 2026-10-12 (Chief): the project own definitions of the reference/slogan/history
% display environments, copied mechanically from the shared Stacks preamble lines 8-12
% (\usepackage{verbatim} plus the three \newenvironment lines, which the chapter preamble keeps
% verbatim from the project preamble). Nothing added, removed or redefined: these environments
% carry no mathematical content and the project itself typesets them as comments. Without these
% three lines a counted statement whose own text uses them fails with Environment slogan undefined.
\usepackage{verbatim}
\newenvironment{reference}{\comment}{\endcomment}
\newenvironment{slogan}{\comment}{\endcomment}
\newenvironment{history}{\comment}{\endcomment}
\usepackage{cleveref}
\crefname{originallemma}{original lemma}{original lemmas}
\Crefname{originallemma}{Original lemma}{Original lemmas}
\makeatletter
\newlabel{algebra-lemma-ui-flat-domain}{{058P}{source}{}{sourcepointer-1}{}}
\newlabel{algebra-lemma-ui-flat-domain@cref}{{[originallemma][0][]058P}{[1][0][]source}{}{}{}}
\makeatother
\crefname{originallemma}{original lemma}{original lemmas}
\Crefname{originallemma}{Original lemma}{Original lemmas}
\makeatletter
\newlabel{lemma-fibres-irreducible-flat-projective-nonnoetherian}{{05FT}{source}{}{sourcepointer-2}{}}
\newlabel{lemma-fibres-irreducible-flat-projective-nonnoetherian@cref}{{[originallemma][0][]05FT}{[1][0][]source}{}{}{}}
\makeatother
\crefname{originallemma}{original lemma}{original lemmas}
\Crefname{originallemma}{Original lemma}{Original lemmas}
\makeatletter
\newlabel{lemma-induction-step-fp}{{05IA}{source}{}{sourcepointer-3}{}}
\newlabel{lemma-induction-step-fp@cref}{{[originallemma][0][]05IA}{[1][0][]source}{}{}{}}
\makeatother
\crefname{originalremark}{original remark}{original remarks}
\Crefname{originalremark}{Original remark}{Original remarks}
\makeatletter
\newlabel{remark-finite-presentation}{{05HY}{source}{}{sourcepointer-4}{}}
\newlabel{remark-finite-presentation@cref}{{[originalremark][0][]05HY}{[1][0][]source}{}{}{}}
\makeatother
\title{A finitely presented flat module with a complete devissage over a prime is, after localising, a direct sum of powers of smooth algebras}
\author{The Stacks Project Authors}
\date{stacks-project:flat.tex:local-snapshot-20261011; GNU FDL 1.2 or later}
\begin{document}
\maketitle
\noindent\textit{Editorial note.} All delivered mathematical text is the author's own, copied from the pinned Stacks Project chapter file flat.tex; no statement or proof step was rewritten, reordered or removed.\par\medskip
\noindent\textit{Editorial note.} Preamble provenance: the chapter's own retained minimal preamble (sources/stacks-flat/retained-preamble.tex, lines 1-85) is a verbatim excerpt of the native Xy-pic setup and of the native theorem declarations and mathematical macros, followed by the project's own reference/slogan/history display-environment definitions copied mechanically from the shared preamble; the wrapper class and the amsmath/amssymb/amsthm/graphicx base packages are supplied by the assembler, and extra\_packages is left empty because everything needed is already in the retained preamble. No macro was invented or redefined.\par\medskip
\noindent\textit{Editorial note.} Same-chapter pointers are located at their \textbackslash{}label lines inside the pinned chapter file; cross-chapter pointers are located in the project's own tag list staged in the snapshot (sources/stacks-chow/tags.txt), at the verbatim lines carrying their canonical tags.\par\medskip
\noindent\textit{Allowed standard prerequisites (editorial).}\begin{itemize}
\item Finiteness of presentation along a complete devissage. Same chapter flat.tex, label remark-finite-presentation, invoked at proof line 3249. Justifies that the first module of the devissage is finitely presented over the first algebra, which is what lets the induction be set up on the first step.
\item The finitely presented induction step. Same chapter flat.tex, label lemma-induction-step-fp, invoked at proof line 3251 and claimed separately as S14-algebra-induction-step-fp. Gives that after inverting an element outside the given prime the first map of the devissage is universally injective over the base, which is the reduction that makes the induction hypothesis applicable.
\item Flatness of the cokernel of a universally injective map over a domain, and projectivity of the smooth pieces. Algebra, label algebra-lemma-ui-flat-domain (proof line 3257); Same chapter flat.tex, label lemma-fibres-irreducible-flat-projective-nonnoetherian (proof line 3268). The first shows that the cokernel of the first map is flat over the base, so that the tail of the devissage applies to it; the second shows that the smooth algebra appearing in the decomposition is projective over the base, which gives the final projectivity statement.
\end{itemize}\medskip
\ifdefined\hypertarget\hypertarget{sourcepointer-1}{}\fi
\noindent\textit{Source reference.} Original lemma 058P: sources/stacks-chow/tags.txt, source lines 6787--6787. Navigation pointer for \textbackslash{}ref\{algebra-lemma-ui-flat-domain\} used inside the counted span (line 3257); the argument that consumes it is retained in the span.\par\medskip
\ifdefined\hypertarget\hypertarget{sourcepointer-2}{}\fi
\noindent\textit{Source reference.} Original lemma 05FT: sources/stacks-flat/pinned-original-source.tex, source lines 2247--2247. Navigation pointer for \textbackslash{}ref\{lemma-fibres-irreducible-flat-projective-nonnoetherian\} used inside the counted span (line 3269); the argument that consumes it is retained in the span.\par\medskip
\ifdefined\hypertarget\hypertarget{sourcepointer-3}{}\fi
\noindent\textit{Source reference.} Original lemma 05IA: sources/stacks-flat/pinned-original-source.tex, source lines 3141--3141. Navigation pointer for \textbackslash{}ref\{lemma-induction-step-fp\} used inside the counted span (line 3251); the argument that consumes it is retained in the span.\par\medskip
\ifdefined\hypertarget\hypertarget{sourcepointer-4}{}\fi
\noindent\textit{Source reference.} Original remark 05HY: sources/stacks-flat/pinned-original-source.tex, source lines 1685--1685. Navigation pointer for \textbackslash{}ref\{remark-finite-presentation\} used inside the counted span (line 3249); the argument that consumes it is retained in the span.\par\medskip
\begin{lemma}
\label{lemma-complete-devissage-flat-finitely-presented-module}
Let $R \to S$ be a ring map of finite presentation.
Let $N$ be a finitely presented $S$-module flat over $R$.
Let $\mathfrak r \subset R$ be a prime ideal.
Assume there exists a complete d\'evissage of $N/S/R$ over $\mathfrak r$.
Then there exists an $f \in R$, $f \not \in \mathfrak r$
such that
$$
N_f \cong B_1^{\oplus r_1} \oplus \ldots \oplus B_n^{\oplus r_n}
$$
as $R$-modules where each $B_i$ is a smooth $R_f$-algebra with geometrically
irreducible fibres. Moreover, $N_f$ is projective as an $R_f$-module.
\end{lemma}

\begin{proof}
Let $(A_i, B_i, M_i, \alpha_i)_{i = 1, \ldots, n}$ be the given
complete d\'evissage. We prove the lemma by induction on $n$.
Note that the assertions of the lemma are entirely about the structure
of $N$ as an $R$-module. Hence we may replace $N$ by $M_1$, and we
may think of $M_1$ as a $B_1$-module. See
Remark \ref{remark-finite-presentation}
in order to see why $M_1$ is of finite presentation as a $B_1$-module. By
Lemma \ref{lemma-induction-step-fp}
we may, after replacing $R$ by $R_f$ for some
$f \in R$, $f \not \in \mathfrak r$, assume
the map $\alpha_1 : B_1^{\oplus r_1} \to M_1$ is $R$-universally injective.
Since $M_1$ and $B_1^{\oplus r_1}$ are $R$-flat and finitely presented as
$B_1$-modules we see that $\Coker(\alpha_1)$ is $R$-flat
(Algebra, Lemma \ref{algebra-lemma-ui-flat-domain})
and finitely presented as a $B_1$-module. Note that
$(A_i, B_i, M_i, \alpha_i)_{i = 2, \ldots, n}$ is a complete
d\'evissage of $\Coker(\alpha_1)$. Hence the induction hypothesis
implies that, after replacing
$R$ by $R_f$ for some $f \in R$, $f \not \in \mathfrak r$,
we may assume that $\Coker(\alpha_1)$ has a decomposition
as in the lemma and is projective. In particular
$M_1 = B_1^{\oplus r_1} \oplus \Coker(\alpha_1)$.
This proves the statement regarding the decomposition.
The statement on projectivity follows as $B_1$ is projective as
an $R$-module by
Lemma \ref{lemma-fibres-irreducible-flat-projective-nonnoetherian}.
\end{proof}

\end{document}
